\begin{center}
    {\LARGE \textbf{2006无机化学}} \\[2ex]
    {\large 时量180分钟 \quad 满分100分}
\end{center}

\vspace{0.5cm}
\noindent\textbf{注意事项:}
\begin{enumerate}[label=\arabic*., parsep=0pt, itemsep=0pt]
    \item 答题前,考生先将自己的姓名、学号、考场号、座位号填写在答题卡上,将条形码准确粘贴在考生信息条形码粘贴区。
    \item 考试结束后,将本试卷与答题纸一并收回。
    \item 回答各个问题时,将答案写在答题纸上,写在本试卷上无效。
\end{enumerate}

\vspace{0.5cm}
\noindent\textbf{一、按要求回答下列问题。第5题10分，其余每题14分。}

\begin{enumerate}[label=\arabic*., start=1, itemsep=1.5ex]
    \item 试写出原子序数为78的元素的名称、符号、电子结构式，并用四个量子数分别表示其最高能级中所有电子，和次高能级中任意两个处于同一轨道中的电子的运动状态。
    \item 试指出下列分子或离子中存在的共轭$\Pi$键（大$\Pi$键）的类型，画图并详细说明它们是怎样形成的。
    \begin{enumerate}[label=(\arabic*), itemsep=1ex]
        \item \ce{CO3^{2-}}
        \item \ce{O3}
    \end{enumerate}
    \item 将反应 \ce{AgCl = Ag^{+} + Cl^{-}} 设计成原电池，试写出两个半反应式。若已知该原电池的电动势为$-0.58\ \text{V}$，且$z=1$，试求\ce{AgCl}的$K_{\text{sp}}$。
    \item 已知下列配离子的分裂能和成对能数据
    \[
    \begin{array}{lcc}
    \hline
     & \ce{[Fe(CN)6]^{4-}} & \ce{[Fe(H2O)6]^{2+}} \\
    \hline
    \Delta/\text{cm}^{-1} & 33000 & 10400 \\
    P/\text{cm}^{-1} & 17600 & 17000 \\
    \hline
    \end{array}
    \]
    \begin{enumerate}[label=(\arabic*), itemsep=1ex]
        \item 试分别用价键理论及晶体场理论解释\ce{[Fe(H2O)6]^{2+}}是高自旋的，\ce{[Fe(CN)6]^{4-}}是低自旋的。
        \item 计算出两种配离子的晶体场稳定化能，并画出必要的简图。
    \end{enumerate}
    \item 完成并配平化学反应方程式
    \begin{enumerate}[label=(\arabic*), itemsep=1ex]
        \item \ce{N2H4 + AgCl ->}
        \item \ce{Co2O3 + HCl ->}
        \item \ce{PbS + H2O2 ->}
        \item \ce{Br2 + Na2CO3 ->}
        \item 金单质溶于王水
    \end{enumerate}

    \item 鉴别
    \begin{enumerate}[label=(\arabic*), itemsep=1ex]
        \item \ce{BeCO3} 和 \ce{MgCO3}
        \item \ce{Na2SO4} 溶液和 \ce{Na2S2O3} 溶液
    \end{enumerate}
    \item 试设计方案分离\ce{Al^{3+}}、\ce{Cr^{3+}}、\ce{Fe^{3+}}、\ce{Co^{2+}}、\ce{Ni^{2+}}共存的离子。
    \item 制备：
    \begin{enumerate}[label=(\arabic*), itemsep=1ex]
        \item 如何配制不含（或极少含有）碳酸钠的氢氧化钠溶液？
        \item 以\ce{NaCl}为基本原料制备\ce{KClO3}，写出反应方程式。
    \end{enumerate}
    \item 有一橙红色固体\ce{A}受热后得绿色的固体\ce{B}和无色的气体\ce{C}，加热时\ce{C}能与镁反应生成灰色的固体\ce{D}。固体\ce{B}溶于过量的\ce{NaOH}溶液生成绿色的溶液\ce{E}，在\ce{E}中加适量\ce{H2O2}则生成黄色溶液\ce{F}。将\ce{F}酸化变为橙色的溶液\ce{G}，在\ce{G}中加\ce{BaCl2}溶液，得黄色沉淀\ce{H}。在\ce{G}中加\ce{KCl}固体，反应完全后则有橙红色晶体\ce{I}析出，滤出\ce{I}烘干并强热则得到的固体产物中有\ce{B}，同时得到能支持燃烧的气体\ce{J}。试写出\ce{A}、\ce{B}、\ce{C}、\ce{D}、\ce{E}、\ce{F}、\ce{G}、\ce{H}、\ce{I}、\ce{J}各物质的化学式，并写出有关的反应方程式。
    \item 试根据下表中所给氮族元素单质的有关数据，讨论氮族元素单质熔、沸点的变化规律并说明原因。
    \[
    \begin{array}{lccccc}
    \hline
    \text{元素} & \ce{N} & \ce{P} & \ce{As} & \ce{Sb} & \ce{Bi} \\
    \hline
    \text{熔点}/^\circ\text{C} & -210.00 & 44.14(\text{白}) & 817(3.91\text{MPa}) & 630.63 & 271.3 \\
    \text{沸点}/^\circ\text{C} & -195.79 & 280.3(\text{白}) & 614(\text{升华}) & 1587 & 1564 \\
    \hline
    \end{array}
    \]
    \item 为什么硫酸盐的热稳定性比碳酸盐高，而硅酸盐的热稳定性更高。
\end{enumerate}